\section{Retornos primitivos y frecuencia levantada}
\subsection{Autómata local de extensión}
Sea \(\mathscr X_\Gamma\) el conjunto finito de estados locales compatibles con la ruta. Un estado conserva

\[
 x=(t,g,B,p,\varepsilon,c,w),
\]
donde \(t\) es el tiempo discreto, \(g\) la fase nonádica, \(B\) el bloque, \(p\) el prefijo, \(\varepsilon\) la hoja, \(c\) el acarreo y \(w\) el devanado. Existe una arista \(x\to y\) cuando \(y\) es una extensión admisible de \(x\), respeta el registro genealógico y supera la poda \(G_9\). La tabla de 36.864 composiciones de acarreo proporciona la ley local de suma; las tablas de retorno fijan la memoria.

Sea \(U_\Gamma\) el operador de transferencia normalizado de este grafo y \(R_\Gamma\) la involución de orientación. Se exige

\[
 U_{C_M\Gamma}=J U_\Gamma J^{-1},\qquad
 R_{C_M\Gamma}=-J R_\Gamma J^{-1}.
\]
\Needspace{7\baselineskip}
\subsection{Conteos primitivos orientados}
Para \(n\ge1\), la traza orientada

\[
 a_n(\Gamma)=\operatorname{Tr}(R_\Gamma U_\Gamma^n)
\]
cuenta retornos cerrados con signo. La inversión de Möbius separa órbitas primitivas:

\[
 p_n(\Gamma)
 =\frac1n\sum_{d\mid n}\mu(d)a_{n/d}(\Gamma).
\]
Se toma el electrón como origen afín de las correcciones relativas y se define

\[
 \boxed{
 \eta_h(\Gamma)=p_h(\Gamma)-p_h(\Gamma_e),
 \qquad h\in\mathfrak S.}
\]
Así se obtiene el mapa

\[
 \mathcal F_{\rm tower}:
 \Gamma^{\rm enr}\longmapsto
 \eta(\Gamma)=(\eta_h(\Gamma))_{h\in\mathfrak S}.
\]
No hay parámetros por especie. El mismo autómata, la misma involución y la misma inversión de Möbius actúan sobre las 324 rutas.

\Needspace{7\baselineskip}
\subsection{Convergencia del refinamiento primitivo}
Sea \(N_\Gamma=\dim\mathscr X_\Gamma\). Si \(U_\Gamma\) es una contracción,

\[
 |a_n(\Gamma)|\le N_\Gamma.
\]
El número de divisores de \(n\) crece sublinealmente, por lo que

\[
 |p_n(\Gamma)|
 \le\frac{N_\Gamma\tau(n)}n.
\]
Como \(0<q_+<q_-<1\),

\[
 |s_h|\le q_-^h+q_+^h\le2q_-^h.
\]
Por comparación,

\[
 \sum_{h\in\mathfrak S}|\eta_h(\Gamma)s_h|<\infty
\]
uniformemente en cualquier colección finita de fibras. El carácter refinado está, por tanto, bien definido.

\Needspace{7\baselineskip}
\subsection{Identidades de control}
La construcción debe satisfacer simultáneamente:

\begin{enumerate}
\item invariancia por renumeración de estados;
\item covarianza bajo conjugación de hoja;
\item compatibilidad con suma directa de componentes;
\item igualdad al recomputar desde registro genealógico y desde el grafo certificado;
\item conservación del origen \(\eta(\Gamma_e)=0\);
\item insensibilidad absoluta a una permutación de las columnas externas de masa;
\item determinación de las colas mediante una cota anterior al contraste.
\end{enumerate}
La representación exacta de un residual dado demuestra completitud del codominio, pero no sustituye estas siete pruebas.

\Needspace{7\baselineskip}
\section{Reloj de ruta y normalización absoluta}
\Needspace{7\baselineskip}
\subsection{Los retornos enteros no bastan para distinguir especies}
En el atlas de 324 rutas, el primer retorno de celda es 27, el retorno cinemático dentro de CT--108 es 54 y el retorno cinemático extendido es 216. Al ser comunes, estos enteros fijan la arquitectura temporal global, pero no pueden explicar por sí solos diferencias de masa entre especies.

\Needspace{7\baselineskip}
\subsection{Holonomía temporal}
La información específica reside en la holonomía del operador de extensión. Sea \(\lambda_j(\Gamma)\) un autovalor no trivial de \(U_\Gamma\) sobre el subespacio persistente. Escribimos

\[
 \lambda_j(\Gamma)=r_j(\Gamma)e^{i\theta_j(\Gamma)}.
\]
La frecuencia interna asociada es

\[
 \omega_j(\Gamma)
 =\frac{\theta_j(\Gamma)+2\pi w_j(\Gamma)}{108t_0},
\]
donde el entero \(w_j\) lo fija el devanado registrado, no una elección de rama posterior. Para un modo estable \(r_j=1\); para una resonancia \(r_j<1\) y su atenuación aporta una anchura.

El cociente

\[
 \rho_{\omega,j}(\Gamma)
 =\frac{\omega_j(\Gamma)}{\omega_0}
\]
es adimensional y puede modificar razones de masa sin alterar la década común \(-34\). La energía y la masa del modo son

\[
 E_j(\Gamma)=\hbar_{\rm ret}\omega_j(\Gamma)
 \chi_{\rm full}(\Gamma)R_{\rm act}^{\kappa_j(\Gamma)},
\]
\[
 m_j(\Gamma)=E_j(\Gamma)c_{\rm int}^{-2}.
\]
Esta ecuación reúne acción, reloj, carácter y lector sin convertir ninguno de ellos en otro. Su aplicación numérica exige determinar \(U_\Gamma\), fijar el devanado y declarar la coordenatización dimensional.
